\subsection{Loading Effect}
\label{sec:theory_general_loading}

Two adjacent circuit stages can load each other if the adjoining
impedances (output -> input or input -> output) are within the same
order of magnitude (say, within a factor of 10 of each other). This
can cause the voltage or current levels to drop significantly as
the adjacent stages form a voltage or current divider. \textbf{Loading
is a pseudo rule and doesn't hold true if under a few special cases
as discussed below}.

For example, Consider three generic circuit stages shown as boxes
connected to each other as shown below. Each circuit stage/box has an
input impedance and an output impedance.
\begin{circuitfig}
  % Paths, nodes and wires:
  \node[shape=rectangle, draw={rgb,255:red,255;green,0;blue,0}, line width=1pt, minimum width=2.465cm, minimum height=1.715cm](N1) at (6.75, 6.625){} node[anchor=center] at (N1.text){\textcolor{rgb,255:red,255;green,0;blue,0}{$Box_1$}};
  \node[shape=rectangle, minimum width=0.965cm, minimum height=0.465cm](N2) at (8.25, 8){} node[anchor=center] at (N2.text){\textcolor{rgb,255:red,255;green,0;blue,0}{$Z_{out_1}$}};
  \node[shape=rectangle, minimum width=0.965cm, minimum height=0.465cm](N3) at (5.5, 8){} node[anchor=center] at (N3.text){\textcolor{rgb,255:red,255;green,0;blue,0}{$Z_{in_1}$}};
  \node[shape=rectangle, draw, line width=1pt, minimum width=2.465cm, minimum height=1.715cm](N4) at (10.75, 6.625){} node[anchor=center] at (N4.text){$Box_2$};
  \node[shape=rectangle, draw={rgb,255:red,0;green,8;blue,255}, line width=1pt, minimum width=2.465cm, minimum height=1.715cm](N5) at (14.5, 6.625){} node[anchor=center] at (N5.text){$Box_3$};
  \draw[stealth-] (9.5, 6.625) -- (8, 6.625);
  \draw[-stealth] (12, 6.625) -- (13.25, 6.625);
  \node[shape=rectangle, minimum width=0.965cm, minimum height=0.465cm](N6) at (12, 8){} node[anchor=center] at (N6.text){\textcolor{rgb,255:red,0;green,0;blue,0}{$Z_{out_2}$}};
  \node[shape=rectangle, minimum width=0.965cm, minimum height=0.465cm](N7) at (15.75, 8){} node[anchor=center] at (N7.text){\textcolor{rgb,255:red,30;green,0;blue,255}{$Z_{out_3}$}};
  \node[shape=rectangle, minimum width=0.965cm, minimum height=0.465cm](N8) at (13.25, 8){} node[anchor=center] at (N8.text){\textcolor{rgb,255:red,0;green,8;blue,255}{$Z_{in_3}$}};
  \node[shape=rectangle, minimum width=0.965cm, minimum height=0.465cm](N9) at (9.5, 8){} node[anchor=center] at (N9.text){$Z_{in_3}$};
  \draw[line width=1pt, dash pattern={on 4pt off 4pt}] (8.75, 4.25) -| (8.75, 9);
  \draw[line width=1pt, dash pattern={on 4pt off 4pt}] (12.5, 4.25) -| (12.5, 9);
\end{circuitfig}
Following our previous definition, If the output impedance of Box 1
(\(Z_{out_1}\)) is negligible compared to the input impedance of Box
2 (\(Z_{in_2}\)), then Box 1 and Box 2 will not load each other
significantly.

Generally, A circuit stage might be a passive circuit, an active circuit,
a mixed (passive and active) circuit, a current source, or an RF circuit.
In these cases, we can summarize the loading effect as follows:

\begin{enumerate}
  \item \textbf{Passive Circuits}: We want to maintain the voltage
    levels without stages loading each other. Therefore, the output
    impedance of the driving stage to be very low (at least 10x lesser)
    than the input impedance of the driven stage. If Box 2 was a Passive
    Circuit, Loading can be minimized if \(Z_{out_2} << 10 x Z_{in_3}\).
  \item \textbf{Active Circuits}: You have to determine what the
    input and output impedance of the specific circuit stage would
    be and then decide if loading is an issue or not between two
    adjacent circuit stages.
  \item \textbf{Mixed Circuits}: Break the circuit down into its
    passive and active components and analyze the loading effect between
    the two adjacent circuit stages.
  \item \textbf{Current Sources}: We want to maintain a constant current
    regardless of the load. To achieve this the output impedance
    of the Current Source should be very high compared to the input
    impedance of the load. Thus, If Box 2 was a Current Source, We would
    want $Z_{out_2}$ to be much larger than $Z_{in_3}$.
  \item \textbf{RF Circuits}: Transmission lines require impedance matching
    to avoid reflections and transfer power from Point A to Point B. In
    this case, if Box 2 and Box 3 were an RF circuit, we would want
    \(Z_{out_2} = Z_{in_3}\). If the impedance were not equal, we place a
    matching network between Box 2 and Box 3 to ensure maximum power transfer.
\end{enumerate}

